%% 附录

\begin{thesisappendix}
\section{六类任务提示词协议}
\label{app:prompt-protocol}

UAV-DualCog 的提示词协议由系统提示词、用户提示模板和结构化输出约束三部分组成。系统提示词统一限定无人机智能体身份、输入媒介、禁止额外自然语言解释以及必须输出可解析结构；用户提示模板注入地标描述、问题文本、候选选项、参考媒体和查询媒体；输出协议则明确绑定答案字段与证据字段。六类任务的协议摘要如\autoref{tab:app-prompt-protocol}所示。

\begin{table}[htbp]
    \centering
    \caption{六类任务提示词协议摘要}
    \label{tab:app-prompt-protocol}
    \small
    \setlength{\tabcolsep}{4pt}
    \begin{tabularx}{\textwidth}{@{}p{2.9cm}p{2.6cm}X>{\raggedright\arraybackslash}p{3.1cm}@{}}
        \toprule
        任务 & 输入媒介 & 核心提示信息 & 必要输出字段 \\
        \midrule
        地标相对位置推理 & 参考图像 + 查询图像 & 地标中心参考视图、当前无人机观察图、候选相对位置选项 & \makecell[l]{\texttt{answer\_option\_id}\\\texttt{bbox\_xyxy\_norm}} \\
        未来观察预测 & 参考图像 + 候选图像 & 当前观察、运动变化描述、多张候选未来观察图 & \makecell[l]{\texttt{answer\_option\_id}\\\texttt{bbox\_xyxy\_norm}} \\
        自我参照位置推理 & 查询图像 & 当前第一视角观察、地标描述、候选方位选项 & \makecell[l]{\texttt{answer\_option\_id}\\\texttt{bbox\_xyxy\_norm}} \\
        地标驱动动作决策 & 查询图像 & 当前观察、目标地标描述、候选动作选项 & \makecell[l]{\texttt{answer\_option\_id}\\\texttt{bbox\_xyxy\_norm}} \\
        飞行行为识别与时间定位 & 视频 & 无人机第一视角飞行视频、复合或原子行为候选项 & \makecell[l]{\texttt{answers}\\\texttt{time\_intervals}} \\
        地标可见次数与区间推理 & 视频 + 参考图像 & 目标地标参考图、第一视角飞行视频、可见性问题 & \makecell[l]{\texttt{count}\\\texttt{time\_intervals}} \\
        \bottomrule
    \end{tabularx}
\end{table}
\end{thesisappendix}

\begin{thesisappendix}
\section{测试基准数据构建实现细节}
\label{app:toolchain-entry}

第四章正文介绍了 UAV-DualCog 的四阶段构建思路。本附录进一步给出对应实现入口、关键伪代码和代表性逻辑，用于说明论文描述与仓库真实实现之间的对应关系。

\begin{table}[htbp]
    \centering
    \caption{测试基准构建工具链主要入口}
    \label{tab:app-toolchain}
    \scriptsize
    \begin{tabularx}{\textwidth}{@{}p{1.7cm}p{3.8cm}X@{}}
        \toprule
        阶段 & 主要入口 & 作用说明 \\
        \midrule
        阶段1 & \makecell[l]{\texttt{stage1\_collect\_}\\\texttt{pcd.py}} & 调用 AerialVLN 模拟器场景扫描，采集语义分割、位姿和点云信息，并融合为场景级几何基础。 \\
        阶段2 & \makecell[l]{\texttt{stage2\_landmark\_}\\\texttt{label.py}} & 基于实例点云生成地标候选，支持人工复核、语义标注和有效地标清单导出。 \\
        阶段3 & \makecell[l]{\texttt{stage3\_generate\_}\\\texttt{traj.py}\\\texttt{stage3\_task\_}\\\texttt{suite.py}} & 生成飞行轨迹、视频资产、复合行为和原子行为时间监督，导出视频任务。 \\
        阶段4 & \makecell[l]{\texttt{stage4\_qa\_generate\_}\\\texttt{and\_eval.py}} & 围绕地标中心参考视图和无人机观察图生成图像问答任务，并导出目标框监督。 \\
        全流程 & \makecell[l]{\texttt{task\_pipeline.py}\\\texttt{pipeline\_common.py}} & 组织阶段间输入输出、路径管理、任务配置和批处理运行。 \\
        \bottomrule
    \end{tabularx}
\end{table}

\subsection{四阶段构建流程伪代码}

\begin{verbatim}
Input: AerialVLN scenes and task configuration
Output: UAV-DualCog benchmark
        with image/video evidence labels

for each scene do
    Stage 1: scan scene and collect RGB / semantics / LiDAR / poses
    fuse observations into scene-level semantic point cloud

    Stage 2: cluster instance candidates from semantic point cloud
    render multi-view landmark candidates
    perform human review and semantic labeling
    export validated landmark assets

    Stage 3: bind landmarks with flight behavior templates
    generate trajectories and videos
    record temporal evidence
    export self-aware and environment-aware video tasks

    Stage 4: sample a landmark-centered reference view
    sample a UAV observation view
    generate self-aware and environment-aware image QA tasks
    compute normalized bounding-box evidence
end for
\end{verbatim}

\subsection{视频任务构建关键逻辑}

阶段3的关键并不只是渲染飞行视频，而是为同一条轨迹同时建立高层行为语义与底层时间证据。其核心流程可以概括为：
\begin{verbatim}
1. select landmark and behavior template
2. generate safe trajectory under flight constraints
3. render first-person UAV video and record timestamps
4. label atomic actions along trajectory
5. aggregate atomic actions into composite spans
6. compute landmark visibility count and visible intervals
7. export structured task manifest for evaluation
\end{verbatim}
这一设计使视频任务能够同时评测“行为识别是否正确”和“时间定位是否可信”。

\subsection{图像任务构建关键逻辑}

阶段4围绕“答案语义”和“空间证据”共同生成图像任务，而不是先生成问题、再事后补框。其核心流程可以概括为：
\begin{verbatim}
1. select target landmark and UAV observation frame
2. build task-specific question and answer options
3. render or retrieve the landmark-centered reference view
4. validate landmark visibility and bounding-box quality
5. compute normalized bounding box in target image
6. export prompt fields, answer option, and bbox evidence
\end{verbatim}
因此，图像任务天生具备“选项判断+目标框定位”的联合监督与联合评测属性。
\end{thesisappendix}

\begin{thesisappendix}
\section{能力优化实现细节}
\label{app:optimization-entry}

第六章中的能力优化部分并非抽象方案，而是已经在仓库中完成实现并通过最新正式训练验证的多模态训练链路。本附录从脚本入口、数据准备伪代码、训练流程、运行完整性核查和关键奖励代码五个角度给出补充说明。

\begin{table}[htbp]
    \centering
    \caption{能力优化主要入口与作用说明}
    \label{tab:app-optimization-entry}
    \scriptsize
    \begin{tabularx}{\textwidth}{@{}p{2.6cm}p{4.1cm}X@{}}
        \toprule
        脚本入口 & 主要输入输出 & 作用说明 \\
        \midrule
        \makecell[l]{\texttt{prepare\_train}\\\texttt{\_dataset.py}} & 输入 Stage 4 manifest；输出 train/valid/all JSONL 与完整性报告 & 将 UAV-DualCog-Train-V1 的图像任务样本转换为统一多模态训练记录，构造训练提示、答案、图像路径与奖励辅助字段，并完成稳定划分。 \\
        \makecell[l]{\texttt{run\_sft}\\\texttt{\_lora.py}} & 输入 train/valid JSONL 与 Qwen 3.5 4B 基座；输出 LoRA adapter、训练指标与合并模型 & 实现第一阶段多模态监督微调，使真实图像通过处理器参与训练，并学习稳定输出结构化 JSON 结果。 \\
        \makecell[l]{\texttt{run\_grpo}\\\texttt{\_light.py}} & 输入训练 JSONL 与 SFT 检查点；输出 GRPO adapter、奖励指标与合并模型 & 实现第二阶段联合 GRPO，在同一奖励函数中同时优化格式、离散答案和目标框质量。 \\
        \makecell[l]{\texttt{run\_full}\\\texttt{\_pipeline.sh}} & 统一管理数据快照、环境检查、日志目录、检查点目录与两阶段调度 & 串联数据准备、SFT 与 Joint GRPO，支持图像分辨率等关键环境参数，形成单机多卡可复现训练流程。 \\
        \makecell[l]{\texttt{check\_training}\\\texttt{\_run.py}} & \makecell[l]{输入运行目录\\输出完整性报告} & 对正式 run 的计划文件、日志、检查点、adapter 与合并模型产物进行完整性核查，并汇总关键训练指标。 \\
        \makecell[l]{\texttt{evaluate}\\\texttt{\_transfer.py}\\\texttt{analyze\_stage}\\\texttt{\_results.py}} & \makecell[l]{输入 Base/最终 GRPO\\benchmark 评测结果；\\SFT 过程结果仅备查} & 负责触发下游 Stage 3/4 评测、聚合跨场景结果，并生成最终优化模型对比表与增益表。 \\
        \bottomrule
    \end{tabularx}
\end{table}

\subsection{训练数据准备伪代码}

\begin{verbatim}
Input: latest stage4 manifests from six training scenes
Output: all / train / valid JSONL and integrity report

for each scene manifest do
    load all image-task samples
    build prompt as benchmark-aligned
        system/user conversational messages
    collect real image paths in benchmark input order
    build assistant completion as JSON:
        {"answer_option_id": "...",
         "bbox_xyxy_norm": [...]}
    cache normalized bbox and bbox_1000
        for reward computation
    record scene-level integrity statistics
end for

bucket samples by (task_type, difficulty)
sort each bucket by sample id
assign every 10th sample to valid split
assign remaining samples to train split
dump JSONL files and integrity report
\end{verbatim}

\subsection{两阶段训练流程伪代码}

\begin{verbatim}
Input: UAV-DualCog-Train-V1 JSONL, Qwen 3.5 4B base model
Output: SFT checkpoint and Joint GRPO checkpoint

Stage 1: SFT
    load train / valid JSONL
    load real RGB images on demand
    encode text and images with AutoProcessor
    train LoRA adapters with completion-only loss
    export adapter and optional merged model

Stage 2: Joint GRPO
    initialize from SFT checkpoint
    load real RGB images on demand
    generate completions under multimodal model
    optimize format reward + answer reward + bbox reward
    export final adapter and optional merged model
\end{verbatim}

\subsection{正式训练完整性核查逻辑}

\begin{verbatim}
Input: training run directory
Output: integrity JSON / Markdown report

check run root exists
for each stage in [01_sft_full, 02_grpo_joint] do
    verify plan file, metrics file, log file exist
    verify adapter and merged_model exist
    count checkpoints and locate latest step
    parse trainer_state tail summary
end for

if both stages complete and merged models exist then
    mark overall_complete = true
\end{verbatim}

\subsection{奖励函数关键代码片段}
\label{app:reward-code}

下列代码片段节选自 \texttt{rewards.py}，展示了本文如何优先解析 JSON 结构化输出，并将“格式合法性、答案正确性、目标框一致性”组合为统一奖励。

\begin{verbatim}
def parse_answer(text: str) -> Optional[str]:
    payload = _extract_json_payload(text)
    if payload is not None:
        answer = str(
            payload.get("answer_option_id", "")
        ).strip().upper()
        if answer in {"A","B","C","D","E","F","G","H"}:
            return answer
    ...

def parse_bbox_norm(text: str) -> Optional[List[float]]:
    payload = _extract_json_payload(text)
    if payload is not None:
        bbox = list(payload.get("bbox_xyxy_norm", []) or [])
        if len(bbox) == 4:
            return _normalize_bbox_scale(bbox)
    ...

def reward_format(text: str) -> float:
    return 1.0 if (parse_answer(text) is not None
                   and parse_bbox_norm(text) is not None) else 0.0

def reward_mcq(pred_text: str, gt_answer: str) -> float:
    pa = parse_answer(pred_text)
    if pa is None:
        return 0.0
    return 1.0 if pa == str(gt_answer).strip().upper() else 0.0

def reward_bbox(pred_text: str, gt_bbox_norm):
    pb = parse_bbox_norm(pred_text)
    if pb is None:
        return 0.0
    g = giou_norm(pb, gt_bbox_norm)
    return max(0.0, min(1.0, (g + 1.0) * 0.5))

def total_reward(pred_text, gt_answer, gt_bbox_norm, weights):
    rf = reward_format(pred_text)
    rm = reward_mcq(pred_text, gt_answer)
    rb = reward_bbox(pred_text, gt_bbox_norm)
    total = (
        weights.w_format * rf
        + weights.w_mcq * rm
        + weights.w_bbox * rb
    )
    return total
\end{verbatim}

这一实现与第六章中的方法公式一一对应，保证奖励设计不是论文中的抽象设想，而是已经进入当前多模态训练链路的可执行机制。
\end{thesisappendix}
